% `tex.web` §§1094, 1126, and 1127: after a source right brace closes an
% ordinary group opened by the template, the entry's brace count is negative.
% The following paragraph inserts another `}`; when that reaches the alignment
% group it is replaced by a frozen `\cr`. With the brace count still negative,
% that boundary in turn inserts `{` before ending the entry and alignment. The
% following explicit `}` must therefore close the surrounding source group.
\catcode`\{=1 \catcode`\}=2 \catcode`\#=6 \catcode`\&=4
\tracingonline=1 \tracingcommands=2 \tracingmacros=1 \tracingrestores=2
\hbadness=10000
\long\def\l#1{#1}
\def\scope{outer}
% A boundary read by `\futurelet` exposes the inaccessible v-template marker.
\halign{x#&x#\cr
  \global\futurelet\endt\foo&\show\endt\cr}
% The source `}` closes the argument begun by the u-template and leaves
% `align_state` negative; the following blank-line `\par` drives all three
% insertions described above.
\noindent{\def\scope{inner}
\halign{\l{#}\cr
  .}

}
\message{PROBE:scope=\scope}
% Far from an active alignment, the same commands are genuinely misplaced;
% unlike their preamble-only sense, `\span` must not expand its successor.
\cr\span
\end
